% GENERATED FILE. Edit canonical lessons, then run npm run generate:books.
% canonical-source-hash: fnv1a64:2b265b3a717bd3c8
% canonical-lessons: HI-C366-vajib, HI-C366-jayaz, HI-C366-bhram, HI-C366-bahkana, HI-C366-majbur, HI-R366-review-359-366

\chapter{Reasonable or Not}
\label{ch:hi-reasonable-or-not}
\hlchaptermodality{366}

\begin{chapteropening}
\textbf{By the end of this chapter:} \emph{I can say whether a demand is fair.}

Complete the last lesson of chapter 366: I can say whether a demand is fair.
\end{chapteropening}

\section[\texorpdfstring{\hi{वाजिब} (vājib)}{वाजिब (vājib)}]{\hi{वाजिब} (vājib) — reasonable, fair (of a price, a demand)}

\label{lesson:HI-C366-vajib}



\begin{quote}
Before the new one: say the Hindi for to hold up, then the Hindi for to get (someone) to accept.
\end{quote}

\subsection*{You'll want to know: \hi{वाजिब}}
\textbf{\hi{वाजिब}} — \emph{vājib} — \textquotedblleft{}reasonable, fair (of a price, a demand)\textquotedblright{}.

\begin{quote}
\hi{यह} \hi{दाम} \hi{वाजिब} \hi{है।}
\end{quote}

\begin{quote}
\emph{yah dām vājib hai.}
\end{quote}

\begin{quote}
This price is reasonable.
\end{quote}

\begin{cousinweb}
\textbf{\hi{वाजिब}} (\emph{vājib}) is usually traced to Arabic \emph{wājib}, \textquotedblleft{}due, necessary\textquotedblright{}, which reached Hindi through Persian.
\end{cousinweb}

\begin{grammarlens}[title={What you've built}]
One more word for giving reasons and opinions.
\end{grammarlens}

\subsection*{Your turn}
\begin{itemize}
\raggedright
  \item \emph{Say it:} \emph{vājib}
  \item \emph{Say it:} \emph{vājib}, once more
  \item \emph{Say it:} \emph{yah dām vājib hai.} — This price is reasonable.
  \item \emph{Recall:} say the Hindi for to hold up, then the Hindi for to get (someone) to accept, then say \emph{vājib} again
\end{itemize}

\begin{tcolorbox}[breakable,colback=teal!4,colframe=teal!35!black,title={Before you move on}]
What does \hi{वाजिब} mean? (\textquotedblleft{}Reasonable, fair (of a price, a demand)\textquotedblright{}.) Say it once more.
\end{tcolorbox}
\section[\texorpdfstring{\hi{जायज़} (jāyaz)}{जायज़ (jāyaz)}]{\hi{जायज़} (jāyaz) — legitimate, justified}

\label{lesson:HI-C366-jayaz}



\begin{quote}
Before the new one: say the Hindi for to get (someone) to accept, then the Hindi for reasonable.
\end{quote}

\subsection*{You'll want to know: \hi{जायज़}}
\textbf{\hi{जायज़}} — \emph{jāyaz} — \textquotedblleft{}legitimate, justified\textquotedblright{}.

\begin{quote}
\hi{तुम्हारा} \hi{ग़ुस्सा} \hi{जायज़} \hi{है।}
\end{quote}

\begin{quote}
\emph{tumhārā ġussā jāyaz hai.}
\end{quote}

\begin{quote}
Your anger is justified.
\end{quote}

\begin{cousinweb}
\textbf{\hi{जायज़}} (\emph{jāyaz}) is usually traced to Arabic \emph{jāʾiz}, \textquotedblleft{}permitted\textquotedblright{}, which reached Hindi through Persian.
\end{cousinweb}

\begin{grammarlens}[title={What you've built}]
One more word for giving reasons and opinions.
\end{grammarlens}

\subsection*{Your turn}
\begin{itemize}
\raggedright
  \item \emph{Say it:} \emph{jāyaz}
  \item \emph{Say it:} \emph{jāyaz}, once more
  \item \emph{Say it:} \emph{tumhārā ġussā jāyaz hai.} — Your anger is justified.
  \item \emph{Recall:} say the Hindi for to get (someone) to accept, then the Hindi for reasonable, then say \emph{jāyaz} again
\end{itemize}

\begin{tcolorbox}[breakable,colback=teal!4,colframe=teal!35!black,title={Before you move on}]
What does \hi{जायज़} mean? (\textquotedblleft{}Legitimate, justified\textquotedblright{}.) Say it once more.
\end{tcolorbox}
\section[\texorpdfstring{\hi{भ्रम} (bhram)}{भ्रम (bhram)}]{\hi{भ्रम} (bhram) — a misconception, an illusion}

\label{lesson:HI-C366-bhram}



\begin{quote}
Before the new one: say the Hindi for reasonable, then the Hindi for legitimate.
\end{quote}

\subsection*{You'll want to know: \hi{भ्रम}}
\textbf{\hi{भ्रम}} — \emph{bhram} — \textquotedblleft{}a misconception, an illusion\textquotedblright{}.

\begin{quote}
\hi{यह} \hi{एक} \hi{भ्रम} \hi{है} \hi{कि} \hi{पैसे} \hi{से} \hi{सब} \hi{कुछ} \hi{मिलता} \hi{है।}
\end{quote}

\begin{quote}
\emph{yah ek bhram hai ki paise se sab kuch miltā hai.}
\end{quote}

\begin{quote}
It's a misconception that money buys everything.
\end{quote}

Masculine.

\begin{cousinweb}
\textbf{\hi{भ्रम}} (\emph{bhram}) is a tatsama, borrowed from Sanskrit \emph{bhrama}, \textquotedblleft{}wandering; confusion\textquotedblright{} (Monier-Williams).
\end{cousinweb}

\begin{grammarlens}[title={What you've built}]
One more word for giving reasons and opinions.
\end{grammarlens}

\subsection*{Your turn}
\begin{itemize}
\raggedright
  \item \emph{Say it:} \emph{bhram}
  \item \emph{Say it:} \emph{bhram}, once more
  \item \emph{Say it:} \emph{yah ek bhram hai ki paise se sab kuch miltā hai.} — It's a misconception that money buys everything.
  \item \emph{Recall:} say the Hindi for reasonable, then the Hindi for legitimate, then say \emph{bhram} again
\end{itemize}

\begin{tcolorbox}[breakable,colback=teal!4,colframe=teal!35!black,title={Before you move on}]
What does \hi{भ्रम} mean? (\textquotedblleft{}A misconception, an illusion\textquotedblright{}.) Say it once more.
\end{tcolorbox}
\section[\texorpdfstring{\hi{बहकाना} (bahkānā)}{बहकाना (bahkānā)}]{\hi{बहकाना} (bahkānā) — to lead astray, to mislead}

\label{lesson:HI-C366-bahkana}



\begin{quote}
Before the new one: say the Hindi for legitimate, then the Hindi for a misconception.
\end{quote}

\subsection*{You'll want to know: \hi{बहकाना}}
\textbf{\hi{बहकाना}} — \emph{bahkānā} — \textquotedblleft{}to lead astray, to mislead\textquotedblright{}.

\begin{quote}
\hi{कुछ} \hi{लोगों} \hi{ने} \hi{बच्चों} \hi{को} \hi{बहकाया।}
\end{quote}

\begin{quote}
\emph{kuch logoṁ ne baccoṁ ko bahkāyā.}
\end{quote}

\begin{quote}
Some people led the children astray.
\end{quote}

Past, with \hi{ने} (he or she …): \hi{उसने} \hi{मुझे} \hi{बहकाया}, \emph{usne mujhe bahkāyā}.

\begin{cousinweb}
\textbf{\hi{बहकाना}} (\emph{bahkānā}) is the causative of \hi{बहकना}, \textquotedblleft{}to be misled\textquotedblright{}, which CDIAL files under Sanskrit \emph{vyathate}, \textquotedblleft{}trembles, goes astray\textquotedblright{} (CDIAL 12164).
\end{cousinweb}

\begin{grammarlens}[title={What you've built}]
One more word for giving reasons and opinions.
\end{grammarlens}

\subsection*{Your turn}
\begin{itemize}
\raggedright
  \item \emph{Say it:} \emph{bahkānā}
  \item \emph{Say it:} \emph{bahkānā}, once more
  \item \emph{Say it:} \emph{kuch logoṁ ne baccoṁ ko bahkāyā.} — Some people led the children astray.
  \item \emph{Recall:} say the Hindi for legitimate, then the Hindi for a misconception, then say \emph{bahkānā} again
\end{itemize}

\begin{tcolorbox}[breakable,colback=teal!4,colframe=teal!35!black,title={Before you move on}]
What does \hi{बहकाना} mean? (\textquotedblleft{}To lead astray, to mislead\textquotedblright{}.) Say it once more.
\end{tcolorbox}
\section[\texorpdfstring{\hi{मजबूर} (majbūr)}{मजबूर (majbūr)}]{\hi{मजबूर} (majbūr) — forced, compelled}

\label{lesson:HI-C366-majbur}



\begin{quote}
Before the new one: say the Hindi for a misconception, then the Hindi for to mislead.
\end{quote}

\subsection*{You'll want to know: \hi{मजबूर}}
\textbf{\hi{मजबूर}} — \emph{majbūr} — \textquotedblleft{}forced, compelled\textquotedblright{}.

\begin{quote}
\hi{बारिश} \hi{की} \hi{वजह} \hi{से} \hi{हम} \hi{घर} \hi{लौटने} \hi{पर} \hi{मजबूर} \hi{हुए।}
\end{quote}

\begin{quote}
\emph{bāriś kī vajah se ham ghar lauṭne par majbūr hue.}
\end{quote}

\begin{quote}
The rain forced us to go back home.
\end{quote}

\begin{cousinweb}
\textbf{\hi{मजबूर}} (\emph{majbūr}) is usually traced to Arabic \emph{majbūr}, \textquotedblleft{}compelled\textquotedblright{}, which reached Hindi through Persian.
\end{cousinweb}

\begin{grammarlens}[title={What you've built}]
One more word for giving reasons and opinions.
\end{grammarlens}

\subsection*{Your turn}
\begin{itemize}
\raggedright
  \item \emph{Say it:} \emph{majbūr}
  \item \emph{Say it:} \emph{majbūr}, once more
  \item \emph{Say it:} \emph{bāriś kī vajah se ham ghar lauṭne par majbūr hue.} — The rain forced us to go back home.
  \item \emph{Recall:} say the Hindi for a misconception, then the Hindi for to mislead, then say \emph{majbūr} again
\end{itemize}

\begin{tcolorbox}[breakable,colback=teal!4,colframe=teal!35!black,title={Before you move on}]
What does \hi{मजबूर} mean? (\textquotedblleft{}Forced, compelled\textquotedblright{}.) Say it once more.
\end{tcolorbox}
\section[(dialogue)]{Looking back: from The Issue at Hand to Reasonable or Not}

\label{lesson:HI-R366-review-359-366}



\begin{quote}
Say the Hindi for an issue and a compromise.
\end{quote}

\subsection*{Your turn}
Say these aloud:

\begin{itemize}
\raggedright
  \item an issue, a compromise, a dispute, to get tangled, to raise (a topic)
  \item an objection, condemnation, a point of view, to threaten, to coax
  \item research, an experiment, a scientist, to gauge, a laboratory
  \item a conclusion, an average, to bring out, an analysis, a trial
  \item a study, according to, a sample, to compare, to assess
  \item an expert, a scholar, an article, an essay, to pick out
  \item a theory, conduct, an exception, to hold up, to get (someone) to accept
  \item reasonable, legitimate, a misconception, to mislead, compelled
\end{itemize}

\begin{tcolorbox}[breakable,colback=teal!4,colframe=teal!35!black,title={Before you move on}]
An issue? (\textbf{\hi{मुद्दा}}, \emph{muddā}.) A compromise? (\textbf{\hi{समझौता}}, \emph{samjhautā}.) A dispute? (\textbf{\hi{विवाद}}, \emph{vivād}.) To get tangled? (\textbf{\hi{उलझना}}, \emph{ulajhnā}.) To raise (a topic)? (\textbf{\hi{छेड़ना}}, \emph{cheṛnā}.) An objection? (\textbf{\hi{आपत्ति}}, \emph{āpatti}.) Condemnation? (\textbf{\hi{निंदा}}, \emph{nindā}.) A point of view? (\textbf{\hi{दृष्टिकोण}}, \emph{dṛṣṭikoṇ}.) To threaten? (\textbf{\hi{धमकाना}}, \emph{dhamkānā}.) To coax? (\textbf{\hi{फुसलाना}}, \emph{phuslānā}.) Research? (\textbf{\hi{शोध}}, \emph{śodh}.) An experiment? (\textbf{\hi{प्रयोग}}, \emph{prayog}.) A scientist? (\textbf{\hi{वैज्ञानिक}}, \emph{vaijñānik}.) To gauge? (\textbf{\hi{मापना}}, \emph{māpnā}.) A laboratory? (\textbf{\hi{प्रयोगशाला}}, \emph{prayogśālā}.) A conclusion? (\textbf{\hi{निष्कर्ष}}, \emph{niṣkarṣ}.) An average? (\textbf{\hi{औसत}}, \emph{ausat}.) To bring out? (\textbf{\hi{उभारना}}, \emph{ubhārnā}.) An analysis? (\textbf{\hi{विश्लेषण}}, \emph{viśleṣaṇ}.) A trial? (\textbf{\hi{परीक्षण}}, \emph{parīkṣaṇ}.) A study? (\textbf{\hi{अध्ययन}}, \emph{adhyayan}.) According to? (\textbf{\hi{अनुसार}}, \emph{anusār}.) A sample? (\textbf{\hi{नमूना}}, \emph{namūnā}.) To compare? (\textbf{\hi{तुलना} \hi{करना}}, \emph{tulnā karnā}.) To assess? (\textbf{\hi{आँकना}}, \emph{ā̃knā}.) An expert? (\textbf{\hi{विशेषज्ञ}}, \emph{viśeṣajña}.) A scholar? (\textbf{\hi{विद्वान}}, \emph{vidvān}.) An article? (\textbf{\hi{लेख}}, \emph{lekh}.) An essay? (\textbf{\hi{निबंध}}, \emph{nibandh}.) To pick out? (\textbf{\hi{छाँटना}}, \emph{chā̃ṭnā}.) A theory? (\textbf{\hi{सिद्धांत}}, \emph{siddhānt}.) Conduct? (\textbf{\hi{व्यवहार}}, \emph{vyavahār}.) An exception? (\textbf{\hi{अपवाद}}, \emph{apvād}.) To hold up? (\textbf{\hi{टिकना}}, \emph{ṭiknā}.) To get (someone) to accept? (\textbf{\hi{मनवाना}}, \emph{manvānā}.) Reasonable? (\textbf{\hi{वाजिब}}, \emph{vājib}.) Legitimate? (\textbf{\hi{जायज़}}, \emph{jāyaz}.) A misconception? (\textbf{\hi{भ्रम}}, \emph{bhram}.) To mislead? (\textbf{\hi{बहकाना}}, \emph{bahkānā}.) Compelled? (\textbf{\hi{मजबूर}}, \emph{majbūr}.)
\end{tcolorbox}
